
\phantomsection
\section*{Заключение}
\addcontentsline{toc}{section}{Заключение}

В результате выполнения курсового проекта была рассмотрена предметная область автоматизации частной медицинской клиники и определены основные проблемы, возникающие при ручной или разрозненной организации записи пациентов, расписания врачей и хранения медицинской информации.

В первой главе проведён анализ существующих программных решений. Сравнение было выполнено в формате «критерий --- платформа»: для Инфоклиника.RU, Medesk, ONDOC, Cliniccards и разрабатываемого приложения рассмотрены онлайн-запись, личный кабинет пациента, электронная медицинская карта, управление расписанием, ролевая модель доступа, работа с документами, администрирование и открытость архитектуры. Анализ показал, что готовые системы обладают широкими возможностями, но являются закрытыми коммерческими продуктами, поэтому для учебного проекта целесообразна разработка собственного веб-прототипа.

Также в первой главе выполнен анализ существующих стеков разработки веб-приложений. Были рассмотрены варианты Python + Flask, Python + Django, Node.js + Express + React/Vue, PHP + Laravel, Java + Spring Boot и ASP.NET Core. По итогам сравнения выбран стек Python + Flask + PostgreSQL + Bootstrap, так как он достаточно функционален для реализации онлайн-записи, медицинских карт, личных кабинетов и CRUD-интерфейсов, но при этом не усложняет архитектуру учебного проекта.

Во второй главе выполнено проектирование веб-приложения. Были определены роли пользователей, описаны пользовательские сценарии, сформулированы функциональные требования, разработаны ER-диаграмма, логическая и физическая схемы базы данных, а также подготовлены макеты страниц. Особое внимание уделено таблицам, связанным с пользователями, расписанием, записями на приём, медицинскими картами и документами.

Таким образом, цель курсового проекта достигнута: выполнены анализ предметной области и проектирование веб-приложения для частной клиники, обеспечивающего онлайн-запись пациентов, работу с расписанием врачей, ведение медицинских карт и предоставление личных кабинетов с разграничением прав доступа. Подготовленные проектные решения могут быть использованы как основа для последующей реализации приложения на стеке Python + Flask. Перспективами развития проекта являются программная реализация модулей авторизации, расписания и медицинских карт, подключение уведомлений по электронной почте или SMS, интеграция с лабораторными сервисами и доработка механизма защиты медицинских данных.
